% !TeX program = xelatex
\documentclass{UAmathtalk}
\renewcommand{\when}{12:00 p.m.}
\renewcommand{\where}{Hudson 0110}
\author{Robin Belton}
\address{Vassar College}
\urladdr{https://robinbelton.github.io/}
\title{Searching for the Unknown: Parameter~Selection~in~Mapper}
\date{Friday, March 26, 2026}


\begin{document}

\maketitle

\begin{abstract}
The Mapper algorithm is a popular tool for visualization and data exploration in TDA with applications ranging from discovering new types of breast cancer to classifying plants. However, it is difficult to use the Mapper algorithm because the user needs to input values for several parameters, and the "best" values are unknown. In this talk, I'll introduce Mapper and how my collaborators, students, and I have approached parameter selection in this setting. In particular, we have focused on studying the inverse problem, proving probabilistic statements, and developing parameter selection algorithms in this setting.
\end{abstract}
\end{document}
